\documentclass[12pt,english]{article}
\usepackage{mathpazo}
\usepackage[T1]{fontenc}
\usepackage[latin9]{inputenc}
\usepackage{float}
\usepackage{amsmath}
\usepackage{amssymb}
\usepackage{geometry}
\geometry{verbose,tmargin=2cm,bmargin=2cm,lmargin=3cm,rmargin=3cm}
\pagestyle{empty}
\usepackage{setspace}
\usepackage[authoryear]{natbib}
\onehalfspacing

\makeatletter
\@ifundefined{date}{}{\date{}}
%%%%%%%%%%%%%%%%%%%%%%%%%%%%%% User specified LaTeX commands.
% general
\usepackage[titletoc]{appendix}
\usepackage{graphicx}
\usepackage{placeins}
\usepackage{tikz}

% algorithm
\usepackage[ruled,vlined,linesnumbered]{algorithm2e}
\IncMargin{2.5cm}
\DecMargin{2cm}
\usepackage{fullwidth}
\usepackage{enumitem}
\setlist{leftmargin=1.7cm}

% tables
\usepackage{tabularx, siunitx, multirow, booktabs}
\begingroup
% Allow `_` and `:` in macro names (LaTeX3 style)
\catcode`\_=11
\catcode`\:=11
% Internal code of `S`
\gdef\tabularxcolumn#1{%
    >{\__siunitx_table_collect_begin:Nn S{} }%
    p{#1}%  <- this is different (is `c` in normal `S`)
    <{\__siunitx_table_print:}%
}
\endgroup

% figures
\usepackage{subfig}
\usepackage{caption}
\captionsetup[subfloat]{position=top}

% footnotes
\setlength{\skip\footins}{1cm}
\usepackage[hang,splitrule]{footmisc}
\setlength{\footnotemargin}{0.3cm} %.5
\setlength{\footnotesep}{0.4cm}

% code
\usepackage{xcolor}
\usepackage{listings}

\definecolor{codegray}{rgb}{0.5,0.5,0.5}
\definecolor{background}{HTML}{F5F5F5}
\definecolor{keyword}{HTML}{4B69C6}
\definecolor{string}{HTML}{448C27}
\definecolor{comment}{HTML}{448C27}

\usepackage{inconsolata}
\lstdefinestyle{mystyle}{
    commentstyle=\color{comment},
    keywordstyle=\color{keyword},
    stringstyle=\color{string},
    basicstyle=\ttfamily,
    breakatwhitespace=false,         
    breaklines=true,                 
    captionpos=b,                    
    keepspaces=true,                                    
    numbersep=5pt,                  
    showspaces=false,                
    showstringspaces=false,
    showtabs=false,
    tabsize=4,
	showlines=true
}

\lstset{style=mystyle}

% manual
\usepackage{enumitem}
\setlist[enumerate]{leftmargin=1cm}
\setlist[itemize]{leftmargin=0.5cm}
% Added by lyx2lyx
\setlength{\parskip}{\smallskipamount}
\setlength{\parindent}{0pt}

\makeatother

\usepackage{babel}
\begin{document}
\title{{\large\textsc{- GEModelTools -}}\\
\textsc{HANC Model}\\
\textsc{with Heterogeneous Returns}}
\author{Rapha\"el Huleux}
\maketitle

\section*{Model}

We consider a \emph{closed} economy with heterogeneous agents and
\emph{flexible prices} \emph{and wages}. Relative to the standard
HANC model, households now differ in the return they earn on their
wealth.

Time is discrete and indexed by $t$. There is a continuum of households
indexed $i$.

\paragraph*{Households.}

Households are heterogeneous \emph{ex post} with respect to their
labor productivity, $e_{it}$, their excess return, $r^{X}_{it}$,
and their assets, $a_{it-1}$. All households share the discount factor
$\beta$. Each period households exogenously supply $\ell_{it}=e_{it}$
units of labor, earn the individual return $r_{t}+r^{X}_{it}$ on
the assets they bring into the period, and choose consumption $c_{it}$
subject to a no-borrowing constraint, $a_{it}\geq0$. Households have
\emph{perfect foresight} wrt. the interest rate and the wage rate,
$\{r_{t},w_{t}\}^{\infty}_{t=0}$, and solve the problem
\begin{align}
v_{t}(e_{it},r^{X}_{it},a_{it-1}) & =\max_{a_{it},c_{it}}\frac{c^{1-\sigma}_{it}}{1-\sigma}+\beta\mathbb{E}_{t}\left[v_{t+1}(e_{it+1},r^{X}_{it+1},a_{it})\right]\label{eq:Bellman}\\
 & \text{s.t.}\nonumber \\
a_{it}+c_{it} & =(1+r_{t}+r^{X}_{it})a_{it-1}+w_{t}e_{it}\nonumber \\
\log e_{it} & =\rho_{e}\log e_{it-1}+\psi_{it},\,\,\,\psi_{it}\sim\mathcal{N}(\mu_{\psi},\sigma^{2}_{\psi}),\,\,\,\mathbb{E}[e_{it}]=1\nonumber \\
r^{X}_{it} & =(1-\rho_{r})\bar{r}^{X}+\rho_{r}r^{X}_{it-1}+\epsilon_{it},\,\,\,\epsilon_{it}\sim\mathcal{N}(0,\sigma^{2}_{r})\nonumber \\
a_{it} & \geq0,\nonumber
\end{align}
where $\psi_{it}$ and $\epsilon_{it}$ are independent, and implicitly
\[
v_{t}(e_{it},r^{X}_{it},a_{it-1})=v(e_{it},r^{X}_{it},a_{it-1},\{r_{\tau},w_{\tau}\}^{\infty}_{\tau=t}).
\]
The excess return $r^{X}_{it}$ is known at the beginning of period
$t$, and has the unconditional mean $\bar{r}^{X}$.

We denote optimal policy functions by $a^{\ast}_{t}(e_{it},r^{X}_{it},a_{it-1})$
and $c^{\ast}_{t}(e_{it},r^{X}_{it},a_{it-1})$. The distribution
of households \emph{before} the realization of idiosyncratic shocks,
i.e. over $e_{it-1}$, $r^{X}_{it-1}$ and $a_{it-1}$, is denoted
$\underline{\boldsymbol{D}}_{t}$. The distribution of households
\emph{after} the realization of idiosyncratic shocks, i.e. over $e_{it}$,
$r^{X}_{it}$ and $a_{it-1}$, is denoted $\boldsymbol{D}_{t}$.

Central aggregate variables are
\begin{align}
A^{hh}_{t} & =\int a^{\ast}_{t}(e_{it},r^{X}_{it},a_{it-1})d\boldsymbol{D}_{t}=\boldsymbol{a}^{\ast\prime}_{t}\boldsymbol{D}_{t}\label{eq:A_hh}\\
C^{hh}_{t} & =\int c^{\ast}_{t}(e_{it},r^{X}_{it},a_{it-1})d\boldsymbol{D}_{t}=\boldsymbol{c}^{\ast\prime}_{t}\boldsymbol{D}_{t}\label{eq:C_hh}\\
RETURNS^{hh}_{t} & =\int(r_{t}+r^{X}_{it})a_{it-1}d\boldsymbol{D}_{t}\label{eq:RETURNS_hh}\\
L^{hh}_{t} & =\int e_{it}d\boldsymbol{D}_{t}=1.\label{eq:L_hh}
\end{align}
$RETURNS^{hh}_{t}$ is the total capital income received by households.
In the standard HANC model it equals $r_{t}A^{hh}_{t-1}$. Here it
does not equal $(r_{t}+\bar{r}^{X})A^{hh}_{t-1}$ in general, because
households with a high $r^{X}_{it}$ accumulate more wealth.

To solve the model, we define the beginning-of-period value function
as
\begin{equation}
\underline{v}_{t}\left(e_{it-1},r^{X}_{it-1},a_{it-1}\right)=\mathbb{E}_{t}\left[v_{t}(e_{it},r^{X}_{it},a_{it-1})\right].\label{eq:beginning-of-period-v}
\end{equation}
The \emph{envelope condition} implies
\begin{equation}
\underline{v}_{t,a}(e_{it-1},r^{X}_{it-1},a_{it-1})=\frac{\partial\underline{v}_{t}(e_{it-1},r^{X}_{it-1},a_{it-1})}{\partial a_{it-1}}=\mathbb{E}_{t}\left[(1+r_{t}+r^{X}_{it})c^{-\sigma}_{it}\right].\label{eq:Envelope}
\end{equation}
Note that the return $r^{X}_{it}$ sits \emph{inside} the expectation.
The \emph{first order condition} for consumption implies
\begin{equation}
c^{-\sigma}_{it}=\beta\underline{v}_{t+1,a}(e_{it},r^{X}_{it},a_{it}).\label{eq:FOC}
\end{equation}

\paragraph*{Firms. }

A representative firm rents capital, $K_{t-1}$, and hires labor, $L_{t}$,
to produce goods, with the production function
\begin{equation}
Y_{t}=\Gamma_{t}K^{\alpha}_{t-1}L^{1-\alpha}_{t},\,\,\,\alpha\in(0,1),\label{eq:Y}
\end{equation}
where $\Gamma_{t}$ is technology. Capital depreciates with the rate
$\delta$. Profit maximization
\[
\max_{K_{t-1},L_{t}}Y_{t}-w_{t}L_{t}-r^{k}_{t}K_{t-1}
\]
implies the standard pricing equations
\begin{align}
r^{k}_{t} & =\alpha\Gamma_{t}(K_{t-1}/L_{t})^{\alpha-1}\label{eq:rk}\\
w_{t} & =(1-\alpha)\Gamma_{t}(K_{t-1}/L_{t})^{\alpha},\label{eq:w}
\end{align}
where $r^{k}_{t}$ is the rental rate of capital and $w_{t}$ is the
wage rate.

\paragraph*{Mutual fund.}

A zero-profit mutual fund owns all the capital. It takes deposits
from households, $A_{t}=K_{t}$, earns $r_{t}=r^{k}_{t}-\delta$ per
unit of capital, and pays household $i$ the return $r_{t}+r^{X}_{it}$.
Zero profit requires that the fund pays out exactly what it earns,
\begin{equation}
RETURNS^{hh}_{t}=r_{t}K_{t-1}.\label{eq:zero_profit}
\end{equation}
Given $A^{hh}_{t-1}=K_{t-1}$, this is the same as $\int r^{X}_{it}a_{it-1}d\boldsymbol{D}_{t}=0$:
the excess returns only redistribute capital income across households,
and the \emph{wealth-weighted} average return equals $r_{t}$.

\paragraph*{Market clearing.}

Market clearing requires
\begin{align}
\text{Capital: }K_{t} & =A_{t}=A^{hh}_{t}\label{eq:clearing_A}\\
\text{Labour: }L_{t} & =L^{hh}_{t}=1\label{eq:clearing_L}\\
\text{Goods: }Y_{t} & =C^{hh}_{t}+\underset{=I_{t}}{\underbrace{K_{t}-K_{t-1}+\delta K_{t-1}}},\label{eq:clearing_Y}
\end{align}
where $I_{t}$ is investment.

\newpage{}

\section*{Steady state and calibration}

We only study the steady state. We normalize $Y=1$ and $L=1$, and
target $K/Y=3$ and $r=0.04$. This gives directly
\begin{align*}
K & =3,\,\,\,\Gamma=\frac{Y}{K^{\alpha}L^{1-\alpha}},\,\,\,r^{k}=\alpha\Gamma(K/L)^{\alpha-1},\,\,\,\delta=r^{k}-r,\,\,\,w=(1-\alpha)\Gamma(K/L)^{\alpha}.
\end{align*}
Given $r$ and $w$, we solve for the two unknowns $\beta$ and $\bar{r}^{X}$
such that
\begin{align*}
A^{hh}-K & =0\,\,\,\text{(capital market clearing)}\\
RETURNS^{hh}-rK & =0\,\,\,\text{(the fund makes zero profit)}.
\end{align*}
Because wealth is concentrated among households with high
$r^{X}_{it}$, the calibrated $\bar{r}^{X}$ is negative.

The remaining parameters are
\begin{center}
\begin{tabular}{llc}
\toprule
Parameter & Description & Value\tabularnewline
\midrule
$\sigma$ & CRRA coefficient & 2.0\tabularnewline
$\alpha$ & Capital share & 0.36\tabularnewline
$\rho_{e}$ & Persistence of $\log e_{it}$ & 0.95\tabularnewline
$\sigma_{\psi}$ & Std. of $\psi_{it}$ & $0.30\sqrt{1-\rho^{2}_{e}}$\tabularnewline
$\rho_{r}$ & Persistence of $r^{X}_{it}$ & 0.8\tabularnewline
$\sigma_{r}$ & Std. of $\epsilon_{it}$ & 0.019\tabularnewline
$N_{e}$ & Grid points for $e_{it}$ & 5\tabularnewline
$N_{r}$ & Grid points for $r^{X}_{it}$ & 7\tabularnewline
\bottomrule
\end{tabular}
\par\end{center}

As a benchmark, we also solve the model without return heterogeneity,
$\sigma_{r}=0$. The calibration then gives $\bar{r}^{X}=0$.

\section*{\newpage Implementation}

The \textbf{files} are: \lstinline{blocks.py}, \lstinline{household_problem.py},
\lstinline{steady_state.py}, and \lstinline{HANCModel.py}.

The results are produced in the \textbf{notebook} \lstinline{HANC.ipynb}.

\textbf{Two exogenous states.} \lstinline{GEModelTools} only has
one stochastic discrete state, $z$. We stack $r^{X}$ and $e$ into
one state with $N_{z}=N_{r}N_{e}$ points. With $\Pi_{r}$ and $\Pi_{e}$
the two transition matrices, the joint transition matrix and ergodic
distribution are the Kronecker products
\[
\Pi_{z}=\Pi_{r}\otimes\Pi_{e},\,\,\,\pi_{z}=\pi_{r}\otimes\pi_{e}.
\]
With $r^{X}$ first, the state \lstinline{i_z} maps to \lstinline{i_r = i_z // Ne}
and \lstinline{i_e = i_z % Ne}.

\textbf{Step 1.} In the \lstinline{.settings} method the household
block now also returns the capital income of each household:

\vspace{2mm}\begin{lstlisting}[language=Python,basicstyle=\linespread{1.3}\ttfamily\footnotesize,
numbers=left,frame=single,backgroundcolor=\color{background}]
self.grids_hh = ['a'] # grids
self.pols_hh = ['a'] # policy functions
self.inputs_hh = ['r','w'] # direct inputs
self.inputs_hh_z = [] # transition matrix inputs
self.outputs_hh = ['a','c','returns'] # outputs
self.intertemps_hh = ['vbeg_a'] # intertemporal variables
\end{lstlisting}\vspace{2mm}

\textbf{Step 2.} In the \lstinline{.setup} method the sizes of the
exogenous grids are:

\vspace{2mm}\begin{lstlisting}[language=Python,basicstyle=\linespread{1.3}\ttfamily\footnotesize,
numbers=left,frame=single,backgroundcolor=\color{background}]
par.Nfix = 1 # number of fixed states
par.Nr = 7 # number of points in r grid
par.Ne = 5 # number of productivity states
par.Nz = par.Ne*par.Nr # number of stochastic discrete states
\end{lstlisting}\vspace{2mm}

\textbf{Step 3.} In the \lstinline{steady_state.py} file the function
\lstinline{prepare_hh_ss} sets the two grids and combines them:

\vspace{2mm}\begin{lstlisting}[language=Python,basicstyle=\linespread{1.3}\ttfamily\footnotesize,
numbers=left,frame=single,backgroundcolor=\color{background}]
par.e_grid[:],e_trans,e_ergodic,_,_ = log_rouwenhorst(par.rho_e,par.sigma_psi,par.Ne)
par.r_grid[:],r_trans,r_ergodic,_,_ = rouwenhorst(par.rX_mean*(1-par.rho_r),par.rho_r,par.sigma_r,par.Nr)
ss.z_trans[0,:,:] = np.kron(r_trans,e_trans) # shape=(Nz,Nz)
ss.Dbeg[0,:,0] = np.kron(r_ergodic,e_ergodic) # all start at a_lag = 0
\end{lstlisting}\vspace{2mm}

\textbf{Step 4.} In the \lstinline{household_problem.py} file the
function \lstinline{solve_hh_backwards} uses the individual return
in the cash-on-hand, in the capital income and in the envelope condition:

\vspace{2mm}\begin{lstlisting}[language=Python,basicstyle=\linespread{1.3}\ttfamily\footnotesize,
numbers=left,frame=single,backgroundcolor=\color{background}]
for i_z in range(par.Nz):
    i_r = i_z // par.Ne
    i_e = i_z % par.Ne
    r_eff = r + par.r_grid[i_r]
    ... # EGM step as in the standard model, with r_eff and e_grid[i_e]
    returns[0,i_z] = r_eff*par.a_grid
    v_a[i_z] = (1+r_eff)*c[0,i_z]**(-par.sigma)
vbeg_a[0] = z_trans[0]@v_a
\end{lstlisting}\vspace{2mm}
\begin{itemize}
\item \vspace{-4mm}\textbf{Note:} $v_{a}$ must be computed for all
$i_{z}$ \emph{before} taking the expectation, since $1+r_{t}+r^{X}_{it}$
differs across states.
\item \vspace{-2mm}\textbf{Check:} Solve the household problem in partial
equilibrium:

\vspace{2mm}\begin{lstlisting}[language=Python,basicstyle=\linespread{1.3}\ttfamily\footnotesize,
numbers=left,frame=single,backgroundcolor=\color{background}]
model.ss.r = 0.04
model.ss.w = 0.64
model.par.beta = 0.89
model.par.rX_mean = 0.0
model.solve_hh_ss(do_print=True) # calls prepare_hh_ss
model.simulate_hh_ss(do_print=True)
model.ss.RETURNS_hh/model.ss.A_hh # average return
\end{lstlisting}
\end{itemize}
\textbf{Step 5.} In the \lstinline{steady_state.py} file the function
\lstinline{find_ss} sets the targets and closed forms, then solves
for $\beta$ and $\bar{r}^{X}$ with \lstinline{scipy.optimize.root},
using the two conditions above:

\vspace{2mm}\begin{lstlisting}[language=Python,basicstyle=\linespread{1.3}\ttfamily\footnotesize,
numbers=left,frame=single,backgroundcolor=\color{background}]
def obj_ss(x,model):
    par.beta = x[0]
    par.rX_mean = x[1]
    model.solve_hh_ss() # calls prepare_hh_ss, which rebuilds the r grid
    model.simulate_hh_ss()
    ss.clearing_A = ss.A-ss.A_hh
    clearing_returns = ss.RETURNS_hh-ss.r*ss.K
    return np.array([ss.clearing_A,clearing_returns])
\end{lstlisting}
\end{document}
